\part{Appearance, labels and evaluation}

\section{Appearance, reusable styles and palettes}
Ten themes are available: blue, teal, solid, gray, lbyellow, viblue, cyblu,
violet, orange and forest. The three new themes pair violet, warm-orange or
forest-green gradients with matching pale tracks and dark value labels.
Nine predefined solid colors are \manualColor{lightgreen}, \manualColor{lightyellow},
\manualColor{lightblue}, \manualColor{lightred}, \manualColor{rose},
\manualColor{skyblue}, \manualColor{gold}, \manualColor{lavender} and
\manualColor{peach}. Use \key{color} for solid fill and
\key{left color}/\key{right color} for a gradient. Standard xcolor expressions work.

\begin{demo}{Violet, warm orange and forest green themes}
\gradsignedbar[min=-50,max=50,label=none,theme=violet]{-24}{24}

\gradsignedbar[min=-50,max=50,label=none,theme=orange]{-24}{24}

\gradsignedbar[min=-50,max=50,label=none,theme=forest]{-24}{24}
\end{demo}


\subsection{Complementary negative colors}
An explicitly selected bar theme also sets negative colors to the RGB complement:
each channel $c$ becomes $1-c$. This covers all ten themes above and
\key{theme=accessible}. Gradient endpoints are complemented individually in the
same left/right order. The solid, gray and accessible themes use solid
complements; gray remains gray with inverted intensity. Without an explicit
theme, the historical solid lightred negative default remains.
\begin{demo}{Built-in signed bars with complementary colors}
\gradsignedbar[min=-50,max=50,error=8,theme=blue]{-24}{24}

\gradsignedbar[min=-50,max=50,error=8,theme=teal]{-24}{24}

\gradsignedbar[min=-50,max=50,error=8,theme=gray]{-24}{24}

\gradsignedbar[min=-50,max=50,theme=blue,
  negative fill=solid,negative color=lightred]{-24}{24}
\end{demo}
This explicitly requested v1.1.1 development change alters negative colors when a
theme is selected; previously most themes retained lightred. Themes reset negative
endpoints and fill mode. Put custom negative settings after the theme. For a solid
override after a gradient theme, set both \key{negative fill=solid} and
\key{negative color}; the latter alone does not replace gradient endpoints.
Later positive endpoint overrides do not recompute negative colors.
Threshold and print precedence, values, errors and scales are unchanged.
The full eleven-theme gallery is in \texttt{docs/examples/theme-complements.tex}.
Complementary colors are not a color-vision accessibility guarantee; retain labels
or other redundant cues.

\begin{demo}{A reusable layout style}
\gradbarsstyle{paperrow}{preset=paper,width=45mm,max=100,precision=0}
\gradbar[style=paperrow]{82}

\graddumbbell[style=paperrow]{60}{82}

\gradbar[preset=report,target=90,thresholds={60,80}]{75}

\gradbar[preset=presentation]{92}

\gradbar[preset=outline]{70}
\end{demo}
The paper preset uses thin solid bars, a faint track and outside values. Report
uses thicker rounded bars; supply actual target and threshold values yourself.
Presentation uses larger type and automatic inside placement. Outline uses hollow
shapes with fine strokes. Presets do not invent data or scale limits.

Group palettes are categorical (six distinct colors), sequential (five blue levels),
diverging (orange through neutral to teal), and mono (gray with textures).
These palettes select by segment position, not by value; short lists cycle.
Named categories instead preserve identity across reordered or subsetted data:
\begin{demo}{Stable category identity and a shared legend}
\gradbarscategory{Compute}{gradbarsBlue}{diagonal}
\gradbarscategory{Storage}{gradbarsOrange}{dots}
\gradstack[stack names={Compute,Storage}]{60,30}

\gradstack[stack names={Storage,Compute}]{30,60}

\gradbarslegend[width=50mm]{Compute,Storage}
\end{demo}
Registered names override positional colors and patterns; unregistered names fall
back to the current palette. Definitions are scoped. In print/mono mode, category
colors use gray while registered patterns remain. Legends use the same mapping.
Patterns are none, diagonal, reverse, dots, crosshatch, horizontal and vertical.
Use names or textures alongside color when category distinctions matter.

\subsection{Color roles with redundant cues}
\key{palette=accessible} supplies seven Okabe--Ito color roles with cyclic stack
textures; \key{theme=accessible} selects its solid-bar roles. Heat provides
\key{heat palette} (\key{accessible-sequential}, \key{accessible-diverging}). Native tables
accept \key{theme=accessible} and \key{style=accessible}. Existing defaults remain.
\begin{demo}{Colors, text, texture and monochrome output}
\gradstack[palette=accessible,width=45mm,height=16pt,
  stack names={A,B,C},segment labels=name]{30,20,50}
\par\medskip
\gradstack[palette=accessible,print,width=45mm,height=16pt,
  stack names={A,B,C},segment labels=name]{30,20,50}
\end{demo}
Colors and textures cycle; keep category text, especially beyond seven categories.
For gradients, center-based automatic contrast is only an estimate; use solid
fills, an external label or a white label background if needed. Heat computes
black/white contrast for the actual cell fill. Missing markers, hollow outlier
circles and printed patterns carry meaning without relying on color alone.
Alternative text retains raw data and explicit units; deliberate \key{alt} can
replace a long description while adjacent text provides the full data.
